\chapter{O Grande Laboratório: Projetos que Podemos Abrir e Verificar}
\label{ch:grande-laboratorio}

Os laboratórios espalhados pelo livro experimentam ideias isoladas. Aqui juntamos peças. Podes começar sem código, avançar para software offline e só depois ligar um modelo real. O objetivo de cada projeto é produzir algo verificável, com uma indicação clara daquilo que continua por demonstrar.

\section{Escolher um percurso e guardar um caderno}
O percurso sem código trabalha com fichas, imagens e decisões. Os exemplos de software dos projetos 3 a 6 usam Python 3.10+ e a biblioteca padrão. Para refazer a exportação da curta com o script do projeto 2, consulta as dependências em \nolinkurl{exemplos/atelier/README.md}. O percurso de inferência local acrescenta modelo, runtime, memória e tempo. Serviços externos são opcionais, com condições e custos a confirmar pelo leitor.

Em cada projeto guarda entrada, objetivo, critérios, versões, saída, avaliação e limitações. O caderno pode ser um ficheiro de texto. Não precisa de parecer um relatório de laboratório espacial para ser útil. Precisa de permitir perceber o que aconteceu e voltar a executar a parte reproduzível.

\begin{promptlabbox}{Projeto 1 — Um museu de afirmações verificadas}
\textbf{Objetivo e contexto:} criar seis cartões que distingam facto, resultado experimental e hipótese sobre IA. Requisitos: papel ou editor de texto, fontes primárias identificadas e, opcionalmente, um chat. Sem código e sem custos obrigatórios.

\textbf{Procedimento:} escolhe duas afirmações de cada tipo; regista a frase, fonte e alcance; pede ao modelo uma explicação acessível; compara a explicação com a ficha; acrescenta uma pergunta que o cartão ainda não responde. Na alternativa offline, escreve a explicação e troca os cartões com outra pessoa.

\textbf{Resultado esperado:} cada cartão tem uma ideia e uma fronteira de generalização. Um resultado sobre quatro semanas não se transforma numa afirmação sobre toda a vida.

\textbf{Validação:} uma pessoa consegue localizar o suporte e dizer o que é observado ou proposto. Regista cartões reprovados e a correção necessária.

\textbf{Falhas possíveis:} inventar referências, resumir apenas títulos ou confundir uma recomendação institucional com um efeito experimental.
\end{promptlabbox}

\begin{promptlabbox}{Projeto 2 — Uma história que se torna cartaz e curta}
\textbf{Objetivo e contexto:} seguir a mesma pista na cena, no cartaz e na curta \emph{A Meia Testemunha}. Podes ver os três materiais em \nolinkurl{exemplos/atelier} e trabalhar com a ilustração fornecida num editor simples.

\textbf{Procedimento:} vê primeiro a curta e depois abre o cartaz. Identifica onde aparece a faixa clara. Faz uma versão que adie a pista e outra que a revele cedo, mantendo os três cortes com cinco segundos cada. Conserva o cartaz com o título numa camada editável e regista a origem da imagem.

\textbf{Resultado esperado:} cada versão conserva a meia, a máquina e o cesto, mas muda a ordem em que reparamos neles. O projeto de referência foi produzido e exportado. Tem cartaz A4, guião visual com três amostras e uma curta muda de quinze segundos.

\textbf{Validação:} confirma os enquadramentos, a duração e os intertítulos no ficheiro. O manifesto \nolinkurl{exemplos/atelier/saidas/resultado.json} regista os resultados da montagem de referência. Para saber como alguém seguiu a pista, usa a pergunta do capítulo~\ref{ch:audiovisual} e anota a observação com a versão vista.

\textbf{Falhas possíveis:} perder a pista num recorte, deixar texto ilegível em ecrã pequeno ou trocar os ficheiros sem atualizar o registo da montagem.
\end{promptlabbox}

\begin{promptlabbox}{Projeto 3 — A Lavandaria encontra uma fronteira de acesso}
\textbf{Objetivo e contexto:} executar o caso integrado já implementado. Requisitos: Python 3.10+, sem dependências externas, rede ou chave; dados inteiramente fictícios.

\textbf{Procedimento:} na raiz, executa \texttt{python3 -m exemplos.lavandaria.avaliar}; consulta LAV-001 e LAV-002 com cliente-A pelo comando do capítulo~\ref{ch:lavandaria}; inspeciona a fonte autorizada e a mensagem de recusa; lê o teste de acesso cruzado.

\textbf{Resultado esperado:} os critérios de regressão conhecidos passam; a consulta autorizada devolve estado; a não autorizada não revela o estado nem confirma existência. A CLI escolhe um papel fictício, sem autenticação real.

\textbf{Validação:} confirma os resultados efetivos e o que o executor recebeu. Este percurso foi executado nesta revisão; o relatório conserva os resultados. O gerador por templates não é um LLM.

\textbf{Falhas possíveis:} testar apenas a recusa e bloquear tudo, confundir uma citação existente com resposta correta ou chamar RAG vetorial à pesquisa lexical.
\end{promptlabbox}

\begin{promptlabbox}{Projeto 4 — Um pacote de evidência que pode ser distribuído}
\textbf{Objetivo e contexto:} rever a skill demonstrativa, sem a instalar automaticamente. Requisitos: Python 3.10+, pacote em \texttt{exemplos/skills/rever-evidencias}; sem rede nem API.

\textbf{Procedimento:} executa \texttt{python3 -m exemplos.laboratorio.skills}; envia o JSON sintético ao script do pacote pelo stdin; cria uma cópia de trabalho e altera um recurso; confirma que a verificação de hash deteta a diferença. Mantém a versão original.

\textbf{Resultado esperado:} pacote original íntegro e fichas com estrutura válida; a cópia alterada exige revisão e nova referência de versão. Ativação num produto é um ensaio separado.

\textbf{Validação:} executa os testes de skills e explica o contraexemplo da afirmação falsa estruturalmente válida. Para distribuir uma adaptação, revê licença e autoria.

\textbf{Falhas possíveis:} atualizar hashes para esconder uma alteração não revista, assumir que um manifesto é uma assinatura ou dar ao pacote permissões que a aplicação não concedeu.
\end{promptlabbox}

\begin{promptlabbox}{Projeto 5 — Uma máquina pequena que aprende e perde precisão}
\textbf{Objetivo e contexto:} ligar treino, avaliação, adaptação e quantização sem descarregar modelos. Requisitos: módulos \texttt{exemplos.laboratorio.fundamentos} e \texttt{exemplos.laboratorio.adaptacao}, Python 3.10+, biblioteca padrão.

\textbf{Procedimento:} executa os dois módulos; identifica treino e teste da reta e do adaptador; verifica a atenção e a tabela de XOR; compara os números reconstruídos pela quantização; duplica o contexto na função de memória. No adaptador, explica por que o segundo alvo conserva erro mesmo depois do treino.

\textbf{Resultado esperado:} a atenção escalar dá cinco, a reta aproxima $2x+1$, a quantização introduz erro e o adaptador conserva a base. O alvo de atualização de posto um é aprendido; no de posto dois, a perda mínima definida no exemplo é $0.25$. A memória ideal mantém hipóteses e exclusões.

\textbf{Validação:} confirma pelo menos uma conta à mão e executa os testes de fundamentos e adaptação. Os dados, configuração e resultados da adaptação estão em \nolinkurl{experiencias/adaptacao_linear}. A implementação foi executada; não treina um Transformer nem mede uma GPU.

\textbf{Falhas possíveis:} transportar o erro minúsculo da reta sem ruído para dados reais ou confundir a grelha didática com um formato industrial de quantização.
\end{promptlabbox}

\begin{promptlabbox}{Projeto 6 — Um incidente numa equipa de especialistas}
\textbf{Objetivo e contexto:} testar orçamento, falha parcial e encaminhamento. Requisitos: Python 3.10+, controlador \texttt{exemplos.laboratorio.equipas}; adaptador stub por omissão.

\textbf{Procedimento:} executa a demonstração; lê os testes; limita o orçamento a duas chamadas; introduz timeout num especialista; tenta devolver um campo extra. Compara estados e contadores. Como extensão, prepara o piloto local da pasta \texttt{experiencias/equipas} com modelo íntegro.

\textbf{Resultado esperado:} o controlador conserva o limite e não apresenta saída parcial como concluída. A extensão com LLM é opt-in e encontra-se bloqueada no registo desta execução.

\textbf{Validação:} os testes devem passar incluindo casos positivos e negativos. Distingue chamadas de função do stub de chamadas reais ao modelo.

\textbf{Falhas possíveis:} ocultar erro numa síntese, repetir ações sem idempotência ou declarar vantagem de multiagentes sem comparação real.
\end{promptlabbox}

\section{Passar de exercício a projeto próprio}
Escolhe uma tarefa pequena e dados que possas usar. Na Lavandaria, a primeira versão pode extrair o assunto e o número da reclamação de uma mensagem. Por exemplo, «Quero saber como vai a LAV-001» tem um assunto e um número. «Ainda espero resposta» tem um assunto possível, mas deixa o número em falta. Para esta tarefa, preencher essa ausência por adivinhação é um erro observável.

\textbf{Primeira passagem em papel.} Escreve três mensagens fictícias. Inclui uma clara, uma com o número em falta e uma com dois assuntos. Ao lado de cada uma, escreve o resultado aceitável, a informação que o sustenta e o que deve ficar por decidir. Guarda uma quarta mensagem para testar a regra depois de a escreveres. Se duas pessoas lerem a regra de maneira diferente, revê a ficha antes de introduzir um modelo.

\textbf{Se quiseres executar.} Usa as mesmas fichas para comparar duas maneiras de resolver a tarefa. Uma pode seguir regras de texto e outra pode pedir a extração a um modelo. Decide antecipadamente que falha obriga a parar, como inventar um número, e regista acerto, tempo e revisão nas mesmas mensagens. O capítulo seguinte mostra como organizar a comparação. Na primeira versão, já tens um resultado se conseguires explicar onde e por que motivo a regra falhou.
